\scanpage{139}
\mspage{88}

et on a
\[
 U_{/t}\simeq U_{/H},
\]
\unsure{et} $U_{/H}$ est interprété en termes du plongement
canonique
\[
 U\hookrightarrow U^\wedge\unsure{.|,}
\]
De \unsure{=|même}, le recollement entre $U'$ et $T'$ sera donné par un préfaisceau
\[
 H'\in\deleted{\unsure{U|U'}}\mathrm{Ob}\,U'^\wedge:\ U'^\circ\unsure{-|\to}\mathrm{Ens},
\]
et on aura
\[
 U'_{/t'}\simeq U_{/H'}.
\]
La relation entre $H'$ et $H$ a été explicitée
au \S~1 ; comme
\[
 H'=p_!H,\quad\text{où}\quad p_!:U^\wedge-U'^\wedge
\]
est l'adjoint à gauche de $p^*:U'^\wedge-U^\wedge$,
donné par
\[
 p^*F'=F'\circ p.
\]
On a
\[
 p_!(H)(u')\simeq\pi_0({}_{u'}\backslash U_{/H}).
\]

\uline{Définition} 1 : \textcircled{a} Soit $U$ un objet de $\mathrm{Cat}$,
$H$ \inserted{$\in\mathrm{Ob}\,U^\wedge$} un préfaisceau sur $U$. On dit que \deleted{\unsure{$F$|le}} $H$
est \uline{totalement} $\mathcal W$-\uline{asphérique} si pour toute
flèche $\unsure{p|f}:U\to U'$ de $\mathrm{Cat}$, $p_!(H)$ de $U'^\wedge$
est un objet $\mathcal W$-asphérique, i.e. $\unsure{W|U|U'}_{/p_!H}$ est
$\mathcal W$-asphérique. \textcircled{b} On dit qu'un objet $V$ de

% archive sheet number: 139
